\section{What the references can and cannot tell apart}\label{sec:results}
\paragraph{Population versus context-only.} Table~\ref{tab:maps} summarises the held-out maps at horizon $h=1$; each map has
its own grid of context lengths, given in the table. With a known noise scale the population prior is worth
\OneAGapMin--\OneAGapMax nats per forecast and is identifiable only at very short contexts. The first study (Stage~0, a
different comparator and rule) found \SzeroAllPairs of \SzeroCells regimes in which all three strategy pairs were
separable. \sHalf, pre-registered with fresh seeds, confirmed a population-specific advantage at one-step but not at
five-step horizons and returned the decision REVISE. With observation noise (F3) the largest $h=1$ gap is
\FthreeGapQzero, \FthreeGapQone, \FthreeGapQtwo and \FthreeGapQthree nats at $q=0,\QOne,\QTwo,\QThree$. No cell is
identifiable at the two larger noise levels or at $h=5$. At moderate noise the F3 map is underpowered: its realised-target
estimator has standard errors several times those of the 1A map, and \FthreeBorderline of its $h=1$ cells (\FthreeBorderlineQone
at $q=\QOne$) are borderline although their point estimates already reach $\delta$. The noise-free F3 row agrees with the
1A map on the three cells compared ($r=1$, $n\le 5$), but the rows differ elsewhere for two reasons visible in the table:
the grids differ (only F3 has $n$ between six and nine, which is why its window at \rsmall ends later), and at $r=4$ the
F3 cells are borderline or unidentifiable where the more precise 1A map is identifiable for the shortest contexts. The
pre-registered \sSix study had already returned STOP for long contexts. When the noise scale is unknown (F1) the window is
much wider: gaps of \FoneGapMin--\FoneGapMax nats, identifiable up to $n=\FoneMaxN$ at every $r$, because the population
also carries information about the noise level, which the context-only reference lacks. In the bivariate family (F2) the
population prior is worth \FtwoGapMin--\FtwoGapMax nats at the longest contexts mapped, identifiable at every tested length
from the second grid point on and at every $r$; the shortest grid point is \textsc{mc-limited}: its gaps are large but the
context-only reference cannot be estimated there to our Monte Carlo criterion.

\begin{table}[t]
\centering\scriptsize
\input{generated/tab_maps.tex}
\caption{Held-out systems, $h=1$: each map's grid of context lengths $n$; the largest $n$ at which population and
context-only Bayes are \textsc{identifiable}, with the range of identifiable gaps in nats; and the number of
\textsc{borderline} cells whose point estimate already reaches $\delta$ (``none'': no identifiable cell; the largest gap
is shown). Generated from the pinned map files.}
\label{tab:maps}
\end{table}

\paragraph{Sensitivity to the margin.} Table~\ref{tab:delta} repeats the held-out comparison at half and twice the
margin. The window moves with $\delta$: with a known noise scale it extends to $n=\FoneMaxN$ at the smaller margin and
shrinks to the two shortest contexts at the larger one. The ordering of the families does not change: the unknown-noise
family keeps the widest window, and the two noisiest settings separate at no margin in this range.

\begin{table}[t]
\centering\small
\input{generated/tab_delta.tex}
\caption{Largest identifiable context length $n$ ($r=1$, $h=1$, held-out) at three margins; ``--'': none. Same files and
three-standard-error rule as Table~\ref{tab:maps}.}
\label{tab:delta}
\end{table}

\paragraph{Library versus continuous prior.} In the scalar families, small libraries are distinguishable from
continuous-prior Bayes on held-out systems and large ones are not; in the bivariate family even large libraries are. In \sHalf, a library of two systems was distinguishable from the continuous prior on
held-out systems in \LibSfiveDist of \LibSfiveCells confirmatory cells. In the 1A map a single-system library is
identifiable in \LibOneAKone of \LibOneACells held-out cells, and every cell with sixteen or more systems is
unidentifiable. In the F1 map every one- and two-system cell is identifiable and \LibFoneKfour of \LibFoneCells
four-system cells are; with sixteen or thirty-two systems every cell is \textsc{partial}, mostly because many draws are
borderline rather than separable. In F3 no held-out cell is identifiable, but for two-system libraries the mean gap reaches
\FthreeLibKtwo nats and the cells are \textsc{partial}. Because a trained model sees one library, Table~\ref{tab:library}
gives the share of individual (cell, library) pairs that are separable: it falls with library size in every scalar family
and stays near one in the bivariate family.

\begin{table}[t]
\centering\small
\input{generated/tab_library.tex}
\caption{Held-out systems, library versus population Bayes: percentage of (cell, library draw) pairs whose lower bound
(estimate $-\,3\,$SE) reaches $\delta$, by library size $K$ (blank: size not mapped). Noisy family at $h=1$. Generated from the
pinned map files.}
\label{tab:library}
\end{table}

On in-library systems the first map found library and continuous posteriors equivalent within $\delta$ in all
\SzeroElibEquiv cells. In the bivariate family the picture changes: on held-out systems a library of sixteen systems is
identifiable in \FtwoLibSixteen of \FtwoLibCells cells and a library of sixty-four in \FtwoLibSixtyFour of
\FtwoLibCells. A plausible reason (not tested here) is that a finite set of systems covers a four-coefficient dynamics
space far more sparsely than an interval of scalar coefficients, so the library posterior stays visibly different from the
continuous prior. In short: in scalar families memorisation of a few systems is visible by score on held-out systems; with a
known noise scale memorisation of sixteen or more is not; with an unknown noise scale separable libraries become rare beyond
eight systems; and in the bivariate family memorisation of sixty-four systems is still visible in most cells.
